%! Author = sriadityavedantam
%! Date = 4/8/26

\section{Chapter 10}

\begin{problem}[Exercise 1]{
    Let $K \leq \zsumz$, and define
    \[
        H \coloneqq \{b \in \z : (a,b) \in K \text{ for some } a \in \z\}.
    \]
    Show that $H \leq \z$.
}
    Given $K \leq \zsumz$, then $(0,0) \in K$.
    Thus by definition of $H$, $0 = e \in H$.
    Given $(a,b) \in K$, then $(a,b)^{-1} = (-a,-b) \in K$.
    Thus if $b \in H$, then $-b \in H$.
    Since $K$ is closed under the operation $(a,b)(c,d) = (a+c,b+d)$, given $b,d \in H$, we have $b+d \in H$.
    Thus $H$ is also closed under the group operation.
    Hence $H \leq \z$.
\end{problem}

\begin{problem}[Exercise 2]{
    Let $G$ be an abelian group, and let $T$ be its torsion subgroup.
}
    \textbf{(a).}
    Show that $T \leq G$.

    Given $G$ is an abelian group and $T$ is the torsion subgroup, then $e \in T$ since the identity has finite order $1$.
    Given $a \in T$, then $a^{-1} \in T$ since the inverse has the same order as $a$.
    Indeed,
    \[
        (a^{-1})^{\ord(a)} = (a^{\ord(a)})^{-1} = e^{-1} = e.
    \]
    Given $a,b \in T$,
    \begin{align*}
        (ab)^{\ord(a)\ord(b)} &= (a^{\ord(a)})^{\ord(b)}(b^{\ord(b)})^{\ord(a)}\\
        &= e^{\ord(b)}e^{\ord(a)}\\
        &= e.
    \end{align*}
    Thus $ab \in T$.
    Hence $T \leq G$.

    \textbf{(b).}
    Show that every nonzero element of $G/T$ has infinite order.

    Given $0 \neq Ta \in G/T$, suppose $a \notin T$ and $Ta$ has finite order $m$.
    Then
    \[
        (Ta)^m = Ta^m = T.
    \]
    Thus $a^m \in T$.
    Since $a^m$ has finite order, there exists $r \in \n$ such that $(a^m)^r = e$.
    Hence $a^{mr} = e$, so $a$ has finite order, and therefore $a \in T$.
    This is a contradiction.
    Hence all nonzero elements of the quotient group $G/T$ have infinite order.
\end{problem}

\begin{problem}[Exercise 3]{
    List the isomorphism types of abelian groups of order $32$.
}
    The isomorphism types of abelian groups are:
    \begin{enumerate}
        \item $2^5$
        \item $2^4 \times 2$
        \item $2^3 \times 2^2$
        \item $2^3 \times 2 \times 2$
        \item $2^2 \times 2^2 \times 2$
        \item $2^2 \times 2 \times 2 \times 2$
        \item $2 \times 2 \times 2 \times 2 \times 2$
    \end{enumerate}
    The representatives for each respectively are:
    \begin{enumerate}
        \item $\z_{32}$
        \item $\z_{16} \times \z_{2}$
        \item $\z_{8} \times \z_{4}$
        \item $\z_{8} \times \z_{2} \times \z_{2}$
        \item $\z_{4} \times \z_{4} \times \z_{2}$
        \item $\z_{4} \times \z_{2} \times \z_{2} \times \z_{2}$
        \item $\z_{2} \times \z_{2} \times \z_{2} \times \z_{2} \times \z_{2}$
    \end{enumerate}
\end{problem}

\begin{problem}[Exercise 4]{
    Rewrite
    \[
        \z_{2} \oplus \z_{2} \oplus \z_{8} \oplus \z_{9} \oplus \z_{5} \oplus \z_{5}
    \]
    in invariant factor form.
}
    \begin{align*}
        2,2,2,8,9,5,5 &\rightarrow 2,2,2,2^3,3^2,5,5\\
        &\implies d_1 = 2^3 \cdot 3^2 \cdot 5 = 360\\
        2,2,2,5 &\implies d_2 = 2 \cdot 5 = 10\\
        2,2 &\implies d_3 = 2.
    \end{align*}
    Thus
    \[
        \mathbb{Z}_2 \oplus \mathbb{Z}_2 \oplus \mathbb{Z}_8 \oplus \mathbb{Z}_9 \oplus \mathbb{Z}_5 \oplus \mathbb{Z}_5 \cong \z_{360} \oplus \z_{10} \oplus \z_{2}.
    \]
\end{problem}

\begin{problem}[Exercise 5]{
    Let $G = A_5$.
}
    \textbf{(a).}
    The possible cycle types are
    \begin{align*}
        5 &\implies (12345),\\
        5 &\implies (12354),\\
        3 + 1 + 1 &\implies (123)(4)(5),\\
        2 + 2 + 1 &\implies (12)(34)(5),\\
        1 + 1 + 1 + 1 + 1 &\implies (1).
    \end{align*}

    \textbf{(b).}
    The number of elements of each type is
    \begin{align*}
        5 &\implies 5!/10 = 12,\\
        5 &\implies 5!/10 = 12,\\
        3 + 1 + 1 &\implies 5 \cdot 4 = 20,\\
        2 + 2 + 1 &\implies 5!/2^3 = 15,\\
        1 + 1 + 1 + 1 + 1 &\implies 1.
    \end{align*}
    Note that we needed to split the $5$-cycles into two conjugacy classes.
    In all summed up, these elements equal $60$, which is the total number of elements of $A_5$.

    \textbf{(c).}
    The size of the centralizer for each element in part \textbf{(a)} is associated with the number of elements above, except for the identity, whose centralizer is all of $G$.

    \textbf{(d).}
    Thus the class equation is
    \[
        \ord(G) = \ord(Z(G)) + [G : C_G(12345)] + [G : C_G(12354)] + [G : C_G(123)] + [G : C_G((12)(34))].
    \]
\end{problem}

\begin{problem}[Exercise 6]{
    Given $K$ is a Sylow $p$-subgroup of $G$ and $N$ is a normal subgroup of $G$, suppose $K \trianglelefteq N$.
    Show that $K \trianglelefteq G$.
}
    Since $K$ is normal in $N$, then $K$ is the only Sylow $p$-subgroup of $N$ by the corollary of the Second Sylow Theorem.
    Let $g \in G$.
    Since $N$ is normal in $G$, we have $gNg^{-1} = N$.
    Thus $gKg^{-1} \leq N$, and $gKg^{-1}$ is also a Sylow $p$-subgroup of $N$.
    Since $K$ is the unique Sylow $p$-subgroup of $N$, it must be that $gKg^{-1} = K$.
    Hence $K \trianglelefteq G$.
\end{problem}

\begin{problem}[Exercise 7]{
    Let $N \trianglelefteq G$, and let $C$ be the conjugacy class of $a$ in $G$.
}
    \textbf{(a).}
    Show that $a \in N$ if and only if $C \subseteq N$.

    \textit{($\implies$).}
    Suppose $a \in N$.
    Since $N$ is normal, for every $x \in G$ we have $xax^{-1} \in N$.
    Thus every conjugate of $a$ lies in $N$, so $C \subseteq N$.

    \textit{($\impliedby$).}
    Suppose $C \subseteq N$.
    Since $a \in C$, it follows immediately that $a \in N$.

    \textbf{(b).}
    Suppose $C_i$ is the conjugacy class of $c \in G$.
    By part \textbf{(a)}, if $c \in N$, then $C_i \subseteq N$.
    By the contrapositive, if $C_i \nsubseteq N$, then $c \notin N$.
    Thus $C_i$ and $N$ are disjoint sets.

    \textbf{(c).}
    The class equation of $G$ is
    \[
        |Z(G)| + \sum_{i=1}^j |C_i|.
    \]
    If we keep only those conjugacy classes contained in $N$, say $C_1,\ldots,C_k$, then by part \textbf{(b)} these classes are either contained in $N$ or disjoint from $N$.
    Thus
    \[
        |N| = \sum_{i=1}^k |C_i|.
    \]
\end{problem}

\begin{problem}[Exercise 8]{
    Given $N \neq \gen{e}$ is a normal subgroup of $G$ and $\ord(G) = p^n$, show that $N \cap Z(G) \neq \gen{e}$.
}
    Since $N$ is a nontrivial $p$-group, its order is $p^k$ for some $k \leq n$.
    By the class equation for $N$, the center $Z(N)$ is nontrivial.
    Since $N$ is normal in $G$, elements of $Z(N)$ lie in $N$.
    Also $Z(N) \subseteq N$, and in particular there exists $e \neq a \in N$ such that $ga g^{-1} = a$ for all $g \in G$.
    Thus $a \in Z(G)$.
    Hence $N \cap Z(G) \neq \gen{e}$.
\end{problem}

\begin{problem}[Exercise 9]{
    Show that $N(N(K)) = N(K)$.
}
    Given $N(K)$ is the normalizer of $K$, we have
    \[
        N(K) = \{x \in G : xKx^{-1} = K\}.
    \]
    Then
    \[
        N(N(K)) = \{x \in G : xN(K)x^{-1} = N(K)\}.
    \]
    Since $N(K)$ is a subgroup, it normalizes itself.
    Hence $N(N(K)) = N(K)$.
\end{problem}

\begin{problem}[Exercise 10]{
    Given $\ord(G) = p^n$, show that $G$ has a normal subgroup of order $p^{n-1}$.
}
    Suppose $k=1$.
    Then a subgroup of order $p^{k-1} = p^0 = 1$ is $\{e\}$.
    Assume that for all $1 \leq i \leq k$ the claim holds, and let $i = k+1$.
    By Theorem 9.27, either $\ord(Z(G)) = p^n$ or $\ord(Z(G)) < p^n$.
    If $\ord(Z(G)) = p^n$, then $G$ is abelian, and the result follows.
    Otherwise, $G/Z(G)$ is a $p$-group of smaller order.
    By induction, $G/Z(G)$ has a normal subgroup of order $p^{n-1-|Z(G)|_p}$, whose inverse image in $G$ is a normal subgroup of order $p^{n-1}$.
    Hence $G$ has a normal subgroup of order $p^{n-1}$.
\end{problem}

\begin{problem}[Exercise 11]{
    Show that there is no simple group of order $30$.
}
    A group of order $30$ has prime divisors $2,3,5$.
    If there is exactly one Sylow $2$-subgroup, then it is normal, and the group cannot be simple.
    Thus there are at least three Sylow $2$-subgroups.

    Similarly, the number of Sylow $3$-subgroups must divide $10$ and be congruent to $1 \mod 3$, so it is at least $10$.

    The number of Sylow $5$-subgroups must divide $6$ and be congruent to $1 \mod 5$, so it is at least $6$.

    Since distinct Sylow $p$-subgroups intersect trivially except at the identity, there are at least
    \[
        3(2-1) + 10(3-1) + 6(5-1) = 3 + 20 + 24 = 47
    \]
    elements of prime order.
    But $47 > 30$, which is impossible.
    Thus a simple group of order $30$ does not exist.
\end{problem}

\begin{problem}[Exercise 12]{
    Given $K$ is a Sylow $p$-subgroup and $N$ is normal in $G$, show that $K \cap N$ is a Sylow $p$-subgroup of $N$.
}
    By Lagrange's Theorem and Exercise 15, we have
    \[
        [N : K \cap N] = [KN : K].
    \]
    Thus
    \[
        [G : K] = [G : KN][KN : K].
    \]
    Hence $[KN : K] \mid [G : K]$.
    Since $[G : K]$ is not divisible by $p$ by the definition of a Sylow $p$-subgroup, it follows that
    \[
        p \nmid [N : K \cap N].
    \]
    Thus $K \cap N$ is a Sylow $p$-subgroup of $N$.
\end{problem}